% Generated by scripts/make_diagrams.py; do not edit by hand.
{\small
\begin{tikzpicture}
  % (a) Type II tree diagram, t channel
  \begin{feynman}
    \vertex (g1) at (0.9, 2.2) {\(g(p_1)\)};
    \vertex (g2) at (0.9, 0.0) {\(g(p_2)\)};
    \vertex[crossed dot] (v1) at (3.2, 2.2) {};
    \vertex[crossed dot] (v2) at (3.2, 0.0) {};
    \vertex (a3) at (5.5, 2.2) {\(A(p_3)\)};
    \vertex (a4) at (5.5, 0.0) {\(A(p_4)\)};
    \diagram* {
      (g1) -- [gluon] (v1) -- [scalar] (a3),
      (g2) -- [gluon] (v2) -- [scalar] (a4),
      (v1) -- [gluon] (v2),
    };
  \end{feynman}
  \node[anchor=north west] at (-0.3, 2.75) {(a)};

  % (b) one-loop Type II diagram (crossed box); loop momentum k
  \begin{scope}[yshift=-3.7cm]
  \begin{feynman}
    \vertex (h1) at (0.9, 2.2) {\(g(p_1)\)};
    \vertex (h2) at (0.9, 0.0) {\(g(p_2)\)};
    \vertex (s1) at (2.0, 2.2);
    \vertex (s2) at (2.0, 0.0);
    \vertex[crossed dot] (e1) at (4.4, 2.2) {};
    \vertex[crossed dot] (e2) at (4.4, 0.0) {};
    \vertex (b1) at (5.7, 2.2) {\(A(p_3)\)};
    \vertex (b2) at (5.7, 0.0) {\(A(p_4)\)};
    \diagram* {
      (h1) -- [gluon] (s1) -- [gluon] (e1) -- [scalar] (b1),
      (h2) -- [gluon] (s2) -- [gluon] (e2) -- [scalar] (b2),
      (s1) -- [gluon] (e2),
      (s2) -- [gluon] (e1),
    };
    % momentum labels placed clear of the gluon coils
    \node at (3.2, 2.62) {\(k\)};
    \node at (3.2, -0.44) {\(p_4-p_1+k\)};
    \node[anchor=east] at (2.45, 1.55) {\(p_1-k\)};
    \node[anchor=west] at (3.95, 1.55) {\(p_3-k\)};
  \end{feynman}
  \node[anchor=north west] at (-0.3, 2.75) {(b)};
  \end{scope}
\end{tikzpicture}}
